\providecommand{\C}{\mathbb{C}}
\providecommand{\Z}{\mathbb{Z}}
\providecommand{\N}{\mathbb{N}}
\providecommand{\bP}{\mathbb{P}}
\providecommand{\cO}{\mathcal{O}}
\providecommand{\Tot}{\mathrm{Tot}}
\providecommand{\ghat}{\widehat{\mathfrak{gl}}_{1}}
\providecommand{\Yplus}{Y^{+}}
\providecommand{\CoHA}{\mathrm{CoHA}}
\providecommand{\Hilb}{\mathrm{Hilb}}
\providecommand{\End}{\mathrm{End}}
\providecommand{\Rep}{\mathrm{Rep}}
\providecommand{\kch}{\kappa_{\mathrm{ch}}}
\providecommand{\kcat}{\kappa_{\mathrm{cat}}}
\providecommand{\tr}{\mathrm{tr}}

\section*{Resolved conifold gluing and the conifold positive half}
\label{sec:resolved-conifold-gluing-positive-half}

Let
\[
  X_{\mathrm{con}}
  =
  \Tot\!\left(\cO_{\bP^1}(-1)^{\oplus 2}\right).
\]
This is one of the two small resolutions of the affine conifold
\[
  Y_0=\operatorname{Spec}\C[x,y,z,w]/(xy-zw).
\]
It is a smooth quasi-projective toric CY$_3$ with one compact curve
$C\simeq\bP^1$, two affine toric charts, and one overlap.  Its base has
dimension one and the bundle has rank two, so it belongs to the
curve-base local CY$_3$ class.
The two charts carry the one-vertex $\C^3$ Hall algebra.  The global
conifold Hall algebra is the two-vertex Klebanov--Witten Jacobi CoHA.
The overlap supplies the descent cocycle and the conifold chamber wall.
Write $\widetilde Y=X_{\mathrm{con}}$ in the toric formulas below.

These are separate layers.  A chart algebra is
$\CoHA(\C^3)\simeq \Yplus(\ghat)$.  The conifold positive half is
$\CoHA(Q_{\mathrm{KW}},W_{\mathrm{KW}})$ before double, centre,
completion, or full current algebra.  The vertex algebra
$\mathcal W_{1+\infty}$ is reached from the Drinfeld double of
$\CoHA(\C^3)$ through the
double/centre/completion/evaluation layer.  The full conifold
Yangian-type object and its $R$-matrix live after Drinfeld doubling, or
after passing to the Drinfeld centre of the $E_1$ representation
category.  At $d=3$ the constructed CY-to-chiral output is native
$E_1$ data; $E_2$-braiding is centre or double data.

\subsection*{The toric fan of $\widetilde{Y}$}

The defining lattice is $N = \Z^{3}$ with rays
\[
  v_{1}=(1,0,0),\quad v_{2}=(0,1,0),\quad v_{3}=(1,0,1),\quad v_{4}=(0,1,1).
\]
These are coplanar on the affine hyperplane
$\langle (1,1,0),v_i\rangle=1$.  This is the toric Calabi--Yau
condition: the anticanonical support function is constant on the rays.
The fan $\Sigma_{\widetilde{Y}}$ has
two maximal cones
\[
  \sigma_{+} = \mathrm{cone}(v_{1},v_{2},v_{3}), \qquad
  \sigma_{-} = \mathrm{cone}(v_{1},v_{2},v_{4}),
\]
sharing the wall $\tau = \mathrm{cone}(v_{1},v_{2})$.  Each
$\sigma_{\pm}$ is generated by a $\Z$-basis of $N$, hence each affine
toric variety $U_{\pm} = \mathrm{Spec}\,\C[\sigma_{\pm}^{\vee}\cap M]$
is $\C^{3}$.  The wall $\tau$ is two-dimensional, so
$U_{+}\cap U_{-} = \mathrm{Spec}\,\C[\tau^{\vee}\cap M] = \C^{*}\times\C^{2}$
with a single $\C^{*}$ direction (the $\bP^{1}$ coordinate away from its
two torus fixed points) and two fibre coordinates.

The distinguished flop $\widetilde{Y}^{\vee}$ is obtained by replacing
$\sigma_{\pm}$ with the other triangulation
$\sigma'_+=\mathrm{cone}(v_{3},v_{4},v_{1})$,
$\sigma'_-=\mathrm{cone}(v_{3},v_{4},v_{2})$;
Bridgeland's theorem
(arXiv:\texttt{math/0502050}) gives $D^{b}(\widetilde{Y})\cong
D^{b}(\widetilde{Y}^{\vee})$ as a derived equivalence.  This is the
flop absorbed by Seiberg duality of the Klebanov--Witten quiver.

\subsection*{Classical gluing at the ring level}

Pick homogeneous $\bP^{1}$-coordinates $[t_{0}:t_{1}]$
and fibre coordinates $u,v$ with weights $(-1,-1)$.  On
$U_{+}=\{t_{0}\neq 0\}$ set $t = t_{1}/t_{0}$ and keep $u,v$ as
$\cO(-1)^{\oplus 2}$ trivialisations; on
$U_{-}=\{t_{1}\neq 0\}$ set $\tilde t = t_{0}/t_{1} = t^{-1}$ and
$\tilde u = t\,u$, $\tilde v = t\,v$ (the $\cO(-1)$ twist absorbed into
the fibre coordinate).  Both charts are $\C^{3}$; the overlap ring is
\[
  \C[t,t^{-1},u,v] \;=\; \C[\tilde t,\tilde t^{-1}, \tilde u,\tilde v]
\]
with cocycle $\tilde t = t^{-1}$, $\tilde u = t u$, $\tilde v = t v$.
The global holomorphic $3$-form $\Omega_{\widetilde{Y}}$ restricts as
\[
  \Omega_{\widetilde{Y}}|_{U_{+}} \;=\; dt\wedge du\wedge dv, \qquad
  \Omega_{\widetilde{Y}}|_{U_{-}} \;=\; -\,d\tilde t\wedge d\tilde u\wedge d\tilde v
\]
because
$d\tilde t\wedge d\tilde u\wedge d\tilde v=-dt\wedge du\wedge dv$.
Non-vanishing of $\Omega_{\widetilde Y}$ on the overlap is the
Calabi--Yau witness.

\subsection*{Chart positive halves and affine-Yangian action}

Each chart has the one-vertex Jordan quiver
$Q_{\C^3}$ with three loops $x,y,z$ and potential
$W_{\C^3}=\tr(x[y,z])$.  Its critical CoHA is the positive Hall half
\[
  \CoHA(Q_{\C^3},W_{\C^3}) \;\simeq\; \Yplus(\ghat)
\]
where the equality is the Schiffmann--Vasserot/Tsymbaliuk
identification of the $\C^3$ critical CoHA with
$Y^{+}(\widehat{\mathfrak{gl}}_{1})$.  The Hilbert-scheme Fock space
$\bigoplus_{n\ge 0}H^*_T(\Hilb^n(\C^2))$ is the standard representation
on which the affine Yangian acts; the chart CoHA supplies its positive
half inside that action.  The charge lattice is
$\Gamma(U_{\pm})=\Z$, with $n\in\N$ labelling the effective length.
For the chart Hall shadow,
\[
  \kch(A_{\C^3}) = 1.
\]
This is a local toric chiral invariant.  The compact categorical Euler
characteristic $\kcat=\chi(\cO_X)$ has compactness hypotheses absent
from an affine chart.

\subsection*{The overlap and the conifold wall}

On $U_{+}\cap U_{-}=\C^{*}\times\C^{2}$ the Laurent transition
\[
  (t,u,v)\longmapsto (t^{-1},tu,tv)
\]
gives the Fourier--Mukai line-bundle cocycle for the two-chart descent.
The same datum appears in the stability space as the conifold wall:
the two simple classes $e_0,e_1$ align, the bound-state class is
$\delta=e_0+e_1$, and the Kontsevich--Soibelman wall automorphism
satisfies the conifold pentagon
\[
  K_{e_0}K_{e_1}=K_{e_1}K_{\delta}K_{e_0}.
\]
In the quantum torus this middle factor is forced:
$x_\delta=q^{-1/2}x_{e_0}x_{e_1}$ when
$\langle e_0,e_1\rangle=1$.  The pentagon fixes the middle wall.
Seiberg duality on the conifold heart exchanges the two real-root
families and fixes $\delta$; geometrically this is the small-resolution
flop.  Morrison--Mozgovoy--Nagao--Szendr\H{o}i
(\texttt{arXiv:1107.5017}) compute this chamber structure and its root
decomposition.

\subsection*{The Klebanov--Witten Jacobi algebra}

The chart quivers are one-vertex Jordan triples.  The global NCCR is
the two-vertex Klebanov--Witten quiver $Q_{\mathrm{KW}}$ with vertices
$0,1$, arrows
$a_{1},a_{2}\colon 0\to 1$ and $b_{1},b_{2}\colon 1\to 0$, and quartic
potential
\[
  W_{\mathrm{KW}} \;=\; \tr(a_{1}b_{1}a_{2}b_{2}-a_{1}b_{2}a_{2}b_{1}).
\]
Its Jacobi algebra is
\[
  J_{\mathrm{KW}}
  =
  \C Q_{\mathrm{KW}}\big/
  \langle \partial_a W_{\mathrm{KW}}: a\in Q_{\mathrm{KW},1}\rangle .
\]
The four cyclic derivatives are
\[
 \partial_{a_1}W_{\mathrm{KW}}=b_1a_2b_2-b_2a_2b_1,\qquad
 \partial_{a_2}W_{\mathrm{KW}}=b_2a_1b_1-b_1a_1b_2,
\]
\[
 \partial_{b_1}W_{\mathrm{KW}}=a_2b_2a_1-a_1b_2a_2,\qquad
 \partial_{b_2}W_{\mathrm{KW}}=a_1b_1a_2-a_2b_1a_1.
\]
Let $e_0,e_1$ be the two vertex idempotents.  The centre is generated
by the central lifts of the four mesons
\[
  \widehat x=a_1b_1+b_1a_1,\qquad
  \widehat y=a_2b_2+b_2a_2,
\]
\[
  \widehat z=a_1b_2+b_2a_1,\qquad
  \widehat w=a_2b_1+b_1a_2.
\]
The $e_0$-corner reads these as
$x=a_1b_1$, $y=a_2b_2$, $z=a_1b_2$, $w=a_2b_1$.  The Jacobi relations
give
\[
  Z(J_{\mathrm{KW}})
  \cong
  \C[\widehat x,\widehat y,\widehat z,\widehat w]/
  (\widehat x\widehat y-\widehat z\widehat w)
  \cong
  \C[x,y,z,w]/(xy-zw).
\]
The algebra is module-finite over this centre and has crepant global
dimension three.  Thus $J_{\mathrm{KW}}$ is the Van den
Bergh--Szendr\H{o}i NCCR of the affine conifold singularity, while the
two small resolutions are its geometric chambers.

\subsection*{Global positive half}

The conifold positive half is the critical CoHA of this global NCCR:
\[
  \Yplus(\widetilde{Y})
  :=
  \CoHA(Q_{\mathrm{KW}},W_{\mathrm{KW}}).
\]
Under the Li--Yamazaki current presentation, with the
Davison--Meinhardt PBW theorem and the MMNS root decomposition as input,
this positive half is
\[
  \CoHA(Q_{\mathrm{KW}},W_{\mathrm{KW}})
  \;\simeq\;
  Y^+\!\bigl(\widehat{\mathfrak{gl}}(1|1)\bigr)^{\mathrm{con}},
\]
as a $\Z_{\geq 0}^2$-graded connected bialgebra in super-vector spaces.
The positive-root pattern is
\[
\Delta_{+}^{\mathrm{re}} = \{(n+1,n),(n,n+1):n\ge 0\},
\qquad
\Delta_{+}^{\mathrm{im}}=\{(n,n):n\ge 1\}.
\]
The real root spaces have superdimension $0|1$ and the imaginary root
spaces have superdimension $2|0$; the two imaginary directions are the
trace and supertrace Cartan modes of the $\widehat{\mathfrak{gl}}(1|1)$
source.  The Drinfeld double
\[
  D\!\left(\Yplus(\widetilde{Y})\right)
  =
  Y\!\bigl(\widehat{\mathfrak{gl}}(1|1)\bigr)^{\mathrm{con}}
\]
is the completed Hopf algebra carrying the universal $R$-matrix.  The
braided $E_2$ layer is equivalently accessed through
$\mathcal Z(\Rep^{E_1}(\Yplus(\widetilde{Y})))$ on a constructed
centre model.

\subsection*{Restriction to charts}

At the level of the MMNS root decomposition and chamber subalgebras,
the open restriction to $U_+\simeq\C^3$ quotients out the classes
supported on the compact curve $C$.  The imaginary roots
$(n,n)$ are curve-supported Cartan modes, hence vanish under this
localisation.  The real roots identify with a single tower indexed by
$n\ge 0$, giving the chart-level positive half $\Yplus(\ghat)$.  The
restriction to $U_-$ gives the same algebra with
$e_n\leftrightarrow f_n$ under Seiberg duality.  The two restrictions
agree on $U_+\cap U_-$ through the Laurent cocycle above.

\subsection*{Invariant normalisation}

The conifold chiral invariant is the direct Jacobi/NCCR value
\[
  \kch(A_{\mathrm{conifold}}) = 1.
\]
Three calculations fix the same normalisation.
\begin{enumerate}
\item The McKay/NCCR calculation uses the tilting bundle
$T=\cO_{X_{\mathrm{con}}}\oplus\pi^*\cO_{\bP^1}(1)$ on
$X_{\mathrm{con}}$, where $\pi:X_{\mathrm{con}}\to\bP^1$ is the bundle
projection, and
$J_{\mathrm{KW}}\cong\End_{X_{\mathrm{con}}}(T)$.  The compactly
supported curve lattice is generated by $[C]$, and the primitive
stable object on $C$ has unit Hall weight.
\item The reduced conifold topological vertex has primitive
Gopakumar--Vafa invariant $n^0_{[C]}=1$; all higher degree curve
contributions are multicovers of this primitive class.
\item The Kontsevich--Soibelman pentagon has a single bound-state wall
$K_\delta$ with exponent one.  The same exponent is the primitive
Hall/chiral contribution.
\end{enumerate}
Thus the conifold row records a direct threefold McKay/NCCR invariant,
separate from compact Hodge supertraces.

The formula
$\kappa_{\mathrm{ch}}(\Tot(\omega_S))=\chi_{\mathrm{top}}(S)/2$ has a
compact Fano surface $S$ as input.  The resolved conifold has base
$\bP^1$ and rank-two bundle $\cO(-1)^{\oplus 2}$, so its input data fall
outside that formula.  The local $\bP^2$ value
$\kappa_{\mathrm{ch}}=3/2$ remains a different toric CY$_3$ example.

The compact invariant $\kcat=\chi(\cO_X)$ requires a finite compact-CY
Euler characteristic.  For the quasi-projective total space
$\widetilde{Y}$, $H^0(\widetilde{Y},\cO_{\widetilde{Y}})$ is
infinite-dimensional.  The Hall/chiral assignment is
$\kch(A_{\mathrm{conifold}})=1$.  Automorphic-product weights and
compact fibre lattices belong to separate constructions.

\subsection*{References}

The NCCR and non-commutative DT model are due to Van den Bergh
(\texttt{arXiv:math/0211064}) and Szendr\H{o}i
(\texttt{arXiv:0705.3419}).  The flop equivalence is Bridgeland's
(\texttt{arXiv:math/0502050}) conifold flop theorem.  The chamber and
root decomposition used above is Morrison--Mozgovoy--Nagao--Szendr\H{o}i
(\texttt{arXiv:1107.5017}).  The PBW input is Davison--Meinhardt
(\texttt{arXiv:1601.02479}); the conifold current presentation is
Li--Yamazaki (\texttt{arXiv:2003.08909}).  The chart identification
$\CoHA(\C^3)\simeq Y^+(\widehat{\mathfrak{gl}}_1)$ is the
Schiffmann--Vasserot/Tsymbaliuk affine-Yangian theorem, with the
Hilbert-scheme Fock action on
$\bigoplus_n H^*_T(\Hilb^n(\C^2))$.
